\section{Lesson 11 --- 29/10/2026}

\subsection{Linear SVM with soft margin}

What if sets $A$ and $B$ are not linearly separable?\\
The linear system
\[
    1 - y^i (w^\top x^i + b) \le 0 \qquad i = 1, \dots, \ell
\]
has no solutions. We add some \g{slack variables} $\xi_i \ge 0$, as few as possible, and we consider the relaxed system
\[
    \begin{aligned}
        1 - y^i (w^\top x^i + b) &\le \xi_i \qquad i = 1, \dots, \ell, \\
        \xi_i &\ge 0 \qquad\ \ \, i = 1, \dots, \ell.
    \end{aligned}
\]
If $x^i$ is misclassified, then $\xi_i > 1$, thus $\sum_{i=1}^{\ell} \xi_i$ is an upper bound of the number of misclassified points. In fact:
\begin{itemize}
    \item $x^i \in A$ with $w^\top x^i + b < 0$ implies $1 < 1 - (w^\top x^i + b) = 1 - y^i (w^\top x^i + b) \le \xi_i$;
    \item $x^i \in B$ with $w^\top x^i + b > 0$ implies $1 < 1 + (w^\top x^i + b) = 1 - y^i (w^\top x^i + b) \le \xi_i$.
\end{itemize}

We redefine the objective function by adding the term $C\sum_{i=1}^{\ell}\xi_i$, where $C > 0$ is a parameter, that represents the error we make:
\begin{equation}
    \begin{cases}
        \displaystyle \min_{w, b, \xi} \ \frac{1}{2} \norm{w}^2 + C \sum_{i=1}^{\ell} \xi_i \\
        1 - y^i (w^\top x^i + b) \le \xi_i \quad \forall i = 1, \dots, \ell, \\
        \xi_i \ge 0 \quad \forall i = 1, \dots, \ell.
    \end{cases}
    \label{eq:svm-soft}
\end{equation}

\begin{example}[Exercise 5.2]
    The dual problem of \cref{eq:svm-soft} is
    \begin{equation}
        \begin{cases}
            \displaystyle \max_{\lambda} \ -\frac{1}{2} \sum_{i=1}^{\ell} \sum_{j=1}^{\ell} y^i y^j (x^i)^\top x^j \lambda_i \lambda_j + \sum_{i=1}^{\ell} \lambda_i \\
            \displaystyle \sum_{i=1}^{\ell} \lambda_i y^i = 0, \\
            0 \le \lambda_i \le C \quad i = 1, \dots, \ell.
        \end{cases}
        \label{eq:svm-soft-dual}
    \end{equation}
    \begin{proof}
        The Lagrangian function is
        \begin{align*}
            L(w, b, \xi, \lambda, \mu) &= \frac{1}{2} \norm{w}^2 + C \sum_{i=1}^{\ell} \xi_i + \sum_{i=1}^{\ell} \lambda_i \bigl(1 - y^i (w^\top x^i + b) - \xi_i\bigr) - \sum_{i=1}^{\ell} \mu_i \xi_i \\
            &= \frac{1}{2} \norm{w}^2 - \sum_{i=1}^{\ell} \lambda_i y^i w^\top x^i - b \sum_{i=1}^{\ell} \lambda_i y^i + \sum_{i=1}^{\ell} \lambda_i + \sum_{i=1}^{\ell} \xi_i (C - \lambda_i - \mu_i).
        \end{align*}
        For fixed $(\lambda, \mu) \in \R^\ell_+ \times \R^\ell_+$, $L$ is convex w.r.t.\ $(w, b, \xi)$, so a global minimum is a stationary point, i.e., a solution of
        \begin{equation}
            \begin{cases}
                \displaystyle \nabla_w L = w - \sum_{i=1}^{\ell} \lambda_i y^i x^i = 0, \\
                \displaystyle \nabla_b L = - \sum_{i=1}^{\ell} \lambda_i y^i = 0, \\
                \nabla_{\xi_i} L = C - \lambda_i - \mu_i = 0 \quad i = 1, \dots, \ell.
            \end{cases}
            \label{eq:soft-stationary}
        \end{equation}
        If \cref{eq:soft-stationary} does not hold, then $\inf L = -\infty$. Otherwise the last sum in $L$ vanishes and, eliminating $w$ with the same arguments used for the dual \cref{eq:svm-dual}, we get
        \[
            \begin{cases}
                \displaystyle \max_{\lambda, \mu} \ -\frac{1}{2} \sum_{i=1}^{\ell} \sum_{j=1}^{\ell} y^i y^j (x^i)^\top x^j \lambda_i \lambda_j + \sum_{i=1}^{\ell} \lambda_i \\
                \displaystyle \sum_{i=1}^{\ell} \lambda_i y^i = 0, \\
                C - \lambda_i - \mu_i = 0 \quad i = 1, \dots, \ell, \\
                \lambda \ge 0, \ \mu \ge 0.
            \end{cases}
        \]
        Eliminating $\mu$, since $\mu_i = C - \lambda_i \ge 0$, we obtain $0 \le \lambda_i \le C$, i.e., \cref{eq:svm-soft-dual}.
    \end{proof}
\end{example}

If $\lambda^*$ is optimal for \cref{eq:svm-soft-dual}, then
\[
    w^* = \sum_{i=1}^{\ell} \lambda^*_i y^i x^i.
\]
We find $b^*$ by choosing $i$ such that $0 < \lambda^*_i < C$ and using the complementarity conditions
\begin{equation}
    \begin{cases}
        \lambda^*_i \bigl(1 - y^i ((w^*)^\top x^i + b^*) - \xi^*_i\bigr) = 0, \\
        \mu^*_i \xi^*_i = (C - \lambda^*_i) \xi^*_i = 0.
    \end{cases}
    \label{eq:soft-compl}
\end{equation}
Thus $b^* = \frac{1}{y^i} - (w^*)^\top x^i$.

\begin{remark}
    From \cref{eq:soft-compl} we have
    \begin{itemize}
        \item $0 \le \lambda^*_i < C$ implies $\xi^*_i = 0$;
        \item $0 < \lambda^*_i \le C$ implies $1 - y^i ((w^*)^\top x^i + b^*) - \xi^*_i = 0$;
        \item $\xi^*_i > 0$ implies $\lambda^*_i = C$.
    \end{itemize}
    These conditions also allow us to find the errors $\xi^*_i$, $i = 1, \dots, \ell$.
\end{remark}

\subsection{Nonlinear SVM}

If two sets $A$ and $B$ are not linearly separable, we ask ourselves: are they linearly separable in another space? To answer this we use a map $\varphi \colon \R^n \to H$, where $H$ is a higher dimensional (maybe infinite) space called \g{feature space}. We try to linearly separate the images $\varphi(x^i)$, $i = 1, \dots, \ell$, in the feature space.

The primal problem is
\begin{equation}
    \begin{cases}
        \displaystyle \min_{w, b, \xi} \ \frac{1}{2} \norm{w}^2 + C \sum_{i=1}^{\ell} \xi_i \\
        1 - y^i (w^\top \varphi(x^i) + b) \le \xi_i \quad \forall i = 1, \dots, \ell, \\
        \xi_i \ge 0 \quad \forall i = 1, \dots, \ell,
    \end{cases}
    \label{eq:svm-nl-primal}
\end{equation}
where $w$ is a vector in a high dimensional space (maybe infinitely many variables).

The dual problem is
\begin{equation}
    \begin{cases}
        \displaystyle \max_{\lambda} \ -\frac{1}{2} \sum_{i=1}^{\ell} \sum_{j=1}^{\ell} y^i y^j \varphi(x^i)^\top \varphi(x^j) \lambda_i \lambda_j + \sum_{i=1}^{\ell} \lambda_i \\
        \displaystyle \sum_{i=1}^{\ell} \lambda_i y^i = 0, \\
        0 \le \lambda_i \le C \quad \forall i = 1, \dots, \ell,
    \end{cases}
    \label{eq:svm-nl-dual}
\end{equation}
and the number of variables is equal to the number of training data.

Let $\lambda^*$ be a solution of the dual problem. Then
\[
    w^* = \sum_{i=1}^{\ell} \lambda^*_i y^i \varphi(x^i).
\]
By any $\lambda^*_i$ such that $0 < \lambda^*_i < C$ we can find $b^*$ by the complementarity relations, that now become
\begin{equation}
    \begin{cases}
        \lambda^*_i \bigl(1 - y^i ((w^*)^\top \varphi(x^i) + b^*) - \xi^*_i\bigr) = 0, \\
        \mu^*_i \xi^*_i = (C - \lambda^*_i) \xi^*_i = 0.
    \end{cases}
    \label{eq:soft-compl-nl}
\end{equation}
Then
\[
    1 - y^i \bigl((w^*)^\top \varphi(x^i) + b^*\bigr) = 0 \quad \longrightarrow \quad b^* = \frac{1}{y^i} - \sum_{j=1}^{\ell} \lambda^*_j y^j \varphi(x^j)^\top \varphi(x^i).
\]
The decision function is given by
\[
    f(x) = \operatorname{sign}\bigl((w^*)^\top \varphi(x) + b^*\bigr) = \operatorname{sign}\Bigl(\sum_{i=1}^{\ell} \lambda^*_i y^i \varphi(x^i)^\top \varphi(x) + b^*\Bigr).
\]
Note that $f(x)$ depends on
\begin{itemize}
    \item $\lambda^*$, that can be found knowing $\varphi(x^i)^\top \varphi(x^j)$;
    \item $\varphi(x^i)^\top \varphi(x)$;
    \item $b^*$, that can be found knowing $\varphi(x^i)^\top \varphi(x^j)$.
\end{itemize}

\subsubsection{Kernel function}

\newpage
\begin{definition}[Kernel function]
    A function $k \colon \R^n \times \R^n \to \R$ is called \g{kernel} if there exists a map $\varphi \colon \R^n \to H$ such that
    \[
        k(x, y) = \langle \varphi(x), \varphi(y) \rangle,
    \]
    where $\langle \cdot, \cdot \rangle$ is a scalar product in the feature space $H$.
\end{definition}

Examples of kernels:
\begin{itemize}
    \item $k(x, y) = x^\top y$;
    \item $k(x, y) = (x^\top y + 1)^p$, with $p \ge 1$ \s{(polynomial)};
    \item $k(x, y) = e^{-\gamma \norm{x - y}^2}$ \s{(Gaussian)};
    \item $k(x, y) = \tanh(\beta x^\top y + \gamma)$, with suitable $\beta$ and $\gamma$.
\end{itemize}

\begin{theorem}
    If $k \colon \R^n \times \R^n \to \R$ is a kernel and $x^1, \dots, x^\ell \in \R^n$, then the matrix $K$ defined by
    \[
        K_{ij} = k(x^i, x^j)
    \]
    is positive semidefinite.
\end{theorem}

The dual problem depends on the kernel $k$:
\[
    \begin{cases}
        \displaystyle \max_{\lambda} \ -\frac{1}{2} \sum_{i=1}^{\ell} \sum_{j=1}^{\ell} y^i y^j k(x^i, x^j) \lambda_i \lambda_j + \sum_{i=1}^{\ell} \lambda_i \\
        \displaystyle \sum_{i=1}^{\ell} \lambda_i y^i = 0, \\
        0 \le \lambda_i \le C \quad i = 1, \dots, \ell.
    \end{cases}
\]

The method is:
\begin{itemize}
    \item choose a kernel $k$;
    \item find an optimal solution $\lambda^*$ of the dual;
    \item choose $i$ such that $0 < \lambda^*_i < C$ and find
    \[
        b^* = \frac{1}{y^i} - \sum_{j=1}^{\ell} \lambda^*_j y^j k(x^i, x^j);
    \]
    \item the decision function is
    \[
        f(x) = \operatorname{sign}\Bigl(\sum_{i=1}^{\ell} \lambda^*_i y^i k(x^i, x) + b^*\Bigr).
    \]
\end{itemize}
The separating surface $f(x) = 0$ is linear in the feature space and nonlinear in the input space.
